\documentclass[12pt]{article}
\usepackage[left = 0.6in, right = 0.6in, top = 0.4in, bottom = 1.5in]{geometry} 
\usepackage{pgfplots, fontawesome5, booktabs, xcolor, hyperref, amsmath,amsthm,amssymb,amsfonts, enumitem, fancyhdr, color, comment, graphicx, environ, actuarialsymbol}
\usepackage[dvipsnames]{xcolor}
\pagestyle{fancy}
\setlength{\headheight}{65pt}
%%%%%%%%%%%%%%%%%%%%%%%%%%%%%%%%%%%%%%%%%%%%%
\lhead{}
\rhead{Trig functions}
%%%%%%%%%%%%%%%%%%%%%%%%%%%%%%%%%%%%%%%%%%%%%
\newcommand{\emptygrid}[3]{ % #1: range, #2: y-label, #3: "labels" or "nolabels"
    \begin{tikzpicture}
    \begin{axis}[
        axis lines = middle,
        xlabel = $x$,
        ylabel = $#2$,
        xmin=-#1, xmax=#1,
        ymin=-#1, ymax=#1,
        grid = both,
        grid style = {line width=.1pt, draw=gray!30},
        xtick distance = 5,
        ytick distance = 5,
        minor tick num = 4,
        tick label style = {font=\tiny},
        width=8cm, height=8cm,
        axis equal image,
        % Logic to hide labels if #3 is "nolabels"
        xticklabels={ \ifx\empty#3\empty\else\ifnum\pdfstrcmp{#3}{nolabels}=0 \else \fi\fi },
        yticklabels={ \ifx\empty#3\empty\else\ifnum\pdfstrcmp{#3}{nolabels}=0 \else \fi\fi }
    ]
    \end{axis}
    \end{tikzpicture}
}
%%%%%%%%%%%%%%%%%%%%%%%%%%%%%%%%%%%%%%%%%%%%%
\newtheorem{theorem}{Theorem}[section]
\begin{document}
\section{Revision}
\subsection{Symmetry properties}

\begin{tikzpicture}[scale=3]
    \draw[->] (-1.2,0) -- (1.2,0) node[right] {$x$};
    \draw[->] (0,-1.2) -- (0,1.2) node[above] {$y$};
    \draw (0,0) circle (1);
\end{tikzpicture}


\begin{itemize}
    \item $\cos(-x) = \cos(x)$
    \item $\sin(-x) = -\sin(x)$
    \item $\tan(-x) = $
    \item $\sin(x+\pi) =$
    \item $\cos(x+\pi) =$
    \item $\sin(x + \frac{\pi}{2}) =$
    \item $\sin(x - \frac{\pi}{2}) =$
    \item $\cos(x + \frac{\pi}{2}) =$
    \item $\cos(x - \frac{\pi}{2}) =$
\end{itemize}

\subsubsection{Pythagorean identity}
\begin{itemize}
    \item $\sin^2(x) + \cos^2(x) = $
    \item $1 + \tan^2(x) = $
    \item $1 + \cot^2(x) = $
\end{itemize}

\subsection{Solving equations}
Strategy:
\begin{enumerate}
    \item Find two distinct solutions in $[0, 2\pi)$ (or $[-\pi, \pi)$)
    \item Use periodicity to find the general solution
\end{enumerate}

Example questions:
\begin{enumerate}
    \item Solve $\sin(2x) = \frac{\sqrt{3}}{2}$
    \item Solve $(\sin(2x) - 1)(\sqrt{3}\sec(\frac{x}{5}) - 2) = 0$
\end{enumerate}

\subsection{Compound angles}
\begin{itemize}
    \item $\sin(a+b) = $
    \item $\cos(a+b) = $
    \item $\tan(a+b) = $
    \item $\sin(a + b) + \sin(a - b) = $
    \item $\cos(a + b) - \cos(a - b) = $
    \item $\tan(2x) = $
\end{itemize}
\begin{enumerate}
    \item let $sec(x) = -3, \quad x \in [\pi ,\frac{3\pi}{2}]$ find the value of $cot(2x)$
    \item Find the value of $tan(\frac{\pi}{8})$
    \item Find without calculus the max and min values of $\sqrt{4sin^{2}(t) + 36cos^{2}(2t)}$
    \item Given that $cosec(x - y) = \frac{5}{4},\quad tan(x)cot(y) = 10$ What is $sin(x+y)?$
\end{enumerate}

\end{document}